\documentclass[10pt, twocolumn, landscape, a4paper]{article}

\usepackage[margin=1in]{geometry}
\usepackage[utf8]{inputenc}
\usepackage[spanish,es-noshorthands]{babel}
\usepackage{amssymb, amsmath, amsbsy}
\usepackage{verbatim}
\usepackage{mathpazo}
\usepackage{tasks}
\usepackage{enumerate}
\usepackage{graphicx}
\usepackage[pdftex]{hyperref}
\hypersetup{pdftitle={conjuntos},pdfauthor={Luis Carrillo Gutiérrez}}

\fontfamily{cmss}\selectfont 
\renewcommand{\familydefault}{\sfdefault}
\pagestyle{empty}

\begin{document}
\subsection*{Determinación de un conjunto}
Si se enumera o nombre cada uno de los elementos de un conjunto ha sido determinado \textbf{por extensión}. Esta forma también se conoce como \emph{FORMA TABULAR} de un conjunto.

\subparagraph*{Ejemplo} --
$$H = \{\text{Lunes}\,;\,\text{Martes}\,;\,\text{Miércoles}\,;\,\text{Jueves}\,;\,\text{Viernes}\,;\,\text{Sábado}\,;\,\text{Domingo}\}$$
$$L = \{a, e, o\}$$
Un conjunto se determina \textbf{por comprensión}, cuando se dá  una o más características o propiedades que cumplan todos y cada uno de los elementos del conjunto. Esta forma también se llama \emph{FORMA CONSTRUCTIVA} de un conjunto.

\subparagraph*{Ejemplo}
$$B = \{x\,/\,x \in A_{3}\}$$
$$D = \{x\,/\,x = 2, 46 < x \leq 93\}$$

\subparagraph*{Observación}
Cuando un conjunto está dado por comprensión, es posible expresarlo por extensión; pero cuando un conjunto esta dado por extensión, \textbf{no} siempre es posible expresarlo por comprensión.
\end{document}